% GENERATED FILE. Edit canonical lessons, then run npm run generate:books.
% canonical-source-hash: fnv1a64:6a4210d8270ea4cf
% canonical-lessons: HI-C304-adhera, HI-C304-garajna, HI-C304-camakna, HI-C304-musladhar, HI-C304-tabahi

\chapter{The Storm}
\label{ch:hi-the-storm}
\hlchaptermodality{304}

\begin{chapteropening}
\textbf{By the end of this chapter:} \emph{I can describe a storm as it came.}

Complete the last lesson of chapter 304: I can describe a storm as it came.
\end{chapteropening}

\section[\texorpdfstring{\hi{अँधेरा} (ãdherā)}{अँधेरा (ãdherā)}]{\hi{अँधेरा} (ãdherā) — darkness}

\label{lesson:HI-C304-adhera}



\begin{quote}
Before the new one: say the Hindi for to cave in, then the Hindi for to be pressed down.
\end{quote}

\subsection*{You'll want to know: \hi{अँधेरा}}
\textbf{\hi{अँधेरा}} — \emph{ãdherā} — \textquotedblleft{}darkness\textquotedblright{}.

\begin{quote}
\hi{दोपहर} \hi{में} \hi{ही} \hi{अँधेरा} \hi{हो} \hi{गया।}
\end{quote}

\begin{quote}
\emph{dopahar meṁ hī ãdherā ho gayā.}
\end{quote}

\begin{quote}
It grew dark in the middle of the afternoon.
\end{quote}

Masculine.

\begin{cousinweb}
\textbf{\hi{अँधेरा}} (\emph{ãdherā}) goes back to Sanskrit \emph{andhakāra}, \textquotedblleft{}darkness\textquotedblright{} (CDIAL 386), built on \emph{andha}, blind. \textbf{\hi{अँधेरा} \hi{होना}} is for it to grow dark; its opposite is \textbf{\hi{उजाला}}.
\end{cousinweb}

\begin{grammarlens}[title={What you've built}]
One more word for telling what happened.
\end{grammarlens}

\subsection*{Your turn}
\begin{itemize}
\raggedright
  \item \emph{Say it:} \emph{ãdherā}
  \item \emph{Say it:} \emph{ãdherā}, once more
  \item \emph{Say it:} \emph{dopahar meṁ hī ãdherā ho gayā.} — It grew dark in the middle of the afternoon.
  \item \emph{Recall:} say the Hindi for to cave in, then the Hindi for to be pressed down, then say \emph{ãdherā} again
\end{itemize}

\begin{tcolorbox}[breakable,colback=teal!4,colframe=teal!35!black,title={Before you move on}]
What does \hi{अँधेरा} mean? (\textquotedblleft{}Darkness\textquotedblright{}.) Say it once more.
\end{tcolorbox}
\section[\texorpdfstring{\hi{गरजना} (garajnā)}{गरजना (garajnā)}]{\hi{गरजना} (garajnā) — to thunder; to roar}

\label{lesson:HI-C304-garajna}



\begin{quote}
Before the new one: say the Hindi for to be pressed down, then the Hindi for darkness.
\end{quote}

\subsection*{You'll want to know: \hi{गरजना}}
\textbf{\hi{गरजना}} — \emph{garajnā} — \textquotedblleft{}to thunder; to roar\textquotedblright{}.

\begin{quote}
\hi{बादल} \hi{ज़ोर} \hi{से} \hi{गरजे।}
\end{quote}

\begin{quote}
\emph{bādal zor se garje.}
\end{quote}

\begin{quote}
The clouds thundered loudly.
\end{quote}

Past (it …): \hi{वह} \hi{शेर} \hi{गरजा}, \emph{vah śer garjā}.

\begin{cousinweb}
\textbf{\hi{गरजना}} (\emph{garajnā}) comes from the Sanskrit verb \emph{garj}, \emph{garjati}, \textquotedblleft{}roars, thunders\textquotedblright{} (Monier-Williams). Clouds \textbf{\hi{गरजते} \hi{हैं}}, and so do lions and angry people.
\end{cousinweb}

\begin{grammarlens}[title={What you've built}]
One more word for telling what happened.
\end{grammarlens}

\subsection*{Your turn}
\begin{itemize}
\raggedright
  \item \emph{Say it:} \emph{garajnā}
  \item \emph{Say it:} \emph{garajnā}, once more
  \item \emph{Say it:} \emph{bādal zor se garje.} — The clouds thundered loudly.
  \item \emph{Recall:} say the Hindi for to be pressed down, then the Hindi for darkness, then say \emph{garajnā} again
\end{itemize}

\begin{tcolorbox}[breakable,colback=teal!4,colframe=teal!35!black,title={Before you move on}]
What does \hi{गरजना} mean? (\textquotedblleft{}To thunder; to roar\textquotedblright{}.) Say it once more.
\end{tcolorbox}
\section[\texorpdfstring{\hi{चमकना} (camaknā)}{चमकना (camaknā)}]{\hi{चमकना} (camaknā) — to flash; to gleam}

\label{lesson:HI-C304-camakna}



\begin{quote}
Before the new one: say the Hindi for darkness, then the Hindi for to thunder.
\end{quote}

\subsection*{You'll want to know: \hi{चमकना}}
\textbf{\hi{चमकना}} — \emph{camaknā} — \textquotedblleft{}to flash; to gleam\textquotedblright{}.

\begin{quote}
\hi{आसमान} \hi{में} \hi{बिजली} \hi{चमकी।}
\end{quote}

\begin{quote}
\emph{āsmān meṁ bijlī camkī.}
\end{quote}

\begin{quote}
Lightning flashed in the sky.
\end{quote}

Past (it …): \hi{वह} \hi{बिजली} \hi{चमकी}, \emph{vah bijlī camkī}.

\begin{cousinweb}
\textbf{\hi{चमकना}} (\emph{camaknā}) is inherited from \emph{*camakka}, a sudden movement, a form CDIAL reconstructs (CDIAL 4676). \textbf{\hi{चमकदार}}, shiny, is built on it.
\end{cousinweb}

\begin{grammarlens}[title={What you've built}]
One more word for telling what happened.
\end{grammarlens}

\subsection*{Your turn}
\begin{itemize}
\raggedright
  \item \emph{Say it:} \emph{camaknā}
  \item \emph{Say it:} \emph{camaknā}, once more
  \item \emph{Say it:} \emph{āsmān meṁ bijlī camkī.} — Lightning flashed in the sky.
  \item \emph{Recall:} say the Hindi for darkness, then the Hindi for to thunder, then say \emph{camaknā} again
\end{itemize}

\begin{tcolorbox}[breakable,colback=teal!4,colframe=teal!35!black,title={Before you move on}]
What does \hi{चमकना} mean? (\textquotedblleft{}To flash; to gleam\textquotedblright{}.) Say it once more.
\end{tcolorbox}
\section[\texorpdfstring{\hi{मूसलाधार} (mūslādhār)}{मूसलाधार (mūslādhār)}]{\hi{मूसलाधार} (mūslādhār) — torrential (of rain)}

\label{lesson:HI-C304-musladhar}



\begin{quote}
Before the new one: say the Hindi for to thunder, then the Hindi for to flash.
\end{quote}

\subsection*{You'll want to know: \hi{मूसलाधार}}
\textbf{\hi{मूसलाधार}} — \emph{mūslādhār} — \textquotedblleft{}torrential (of rain)\textquotedblright{}.

\begin{quote}
\hi{शाम} \hi{को} \hi{मूसलाधार} \hi{बारिश} \hi{हुई।}
\end{quote}

\begin{quote}
\emph{śām ko mūslādhār bāriś huī.}
\end{quote}

\begin{quote}
In the evening there was torrential rain.
\end{quote}

\begin{cousinweb}
\textbf{\hi{मूसलाधार}} (\emph{mūslādhār}) joins Sanskrit \emph{musala}, \textquotedblleft{}a pestle\textquotedblright{}, and \emph{dhārā}, \textquotedblleft{}a stream\textquotedblright{} (Monier-Williams): rain falling in streams as thick as a pestle. It is used of rain and of nothing else.
\end{cousinweb}

\begin{grammarlens}[title={What you've built}]
One more word for telling what happened.
\end{grammarlens}

\subsection*{Your turn}
\begin{itemize}
\raggedright
  \item \emph{Say it:} \emph{mūslādhār}
  \item \emph{Say it:} \emph{mūslādhār}, once more
  \item \emph{Say it:} \emph{śām ko mūslādhār bāriś huī.} — In the evening there was torrential rain.
  \item \emph{Recall:} say the Hindi for to thunder, then the Hindi for to flash, then say \emph{mūslādhār} again
\end{itemize}

\begin{tcolorbox}[breakable,colback=teal!4,colframe=teal!35!black,title={Before you move on}]
What does \hi{मूसलाधार} mean? (\textquotedblleft{}Torrential (of rain)\textquotedblright{}.) Say it once more.
\end{tcolorbox}
\section[\texorpdfstring{\hi{तबाही} (tabāhī)}{तबाही (tabāhī)}]{\hi{तबाही} (tabāhī) — destruction; havoc}

\label{lesson:HI-C304-tabahi}



\begin{quote}
Before the new one: say the Hindi for to flash, then the Hindi for torrential (of rain).
\end{quote}

\subsection*{You'll want to know: \hi{तबाही}}
\textbf{\hi{तबाही}} — \emph{tabāhī} — \textquotedblleft{}destruction; havoc\textquotedblright{}.

\begin{quote}
\hi{तूफ़ान} \hi{ने} \hi{पूरे} \hi{शहर} \hi{में} \hi{तबाही} \hi{की।}
\end{quote}

\begin{quote}
\emph{tūfān ne pūre śahar meṁ tabāhī kī.}
\end{quote}

\begin{quote}
The storm wreaked havoc across the whole city.
\end{quote}

Feminine.

\begin{cousinweb}
\textbf{\hi{तबाही}} (\emph{tabāhī}) is usually traced to Persian \emph{tabāhī}, \textquotedblleft{}ruin\textquotedblright{}, from \emph{tabāh}, ruined. A storm or a war \textbf{\hi{तबाही} \hi{करता} \hi{है}}, or \textbf{\hi{तबाही} \hi{मचाता} \hi{है}}.
\end{cousinweb}

\begin{grammarlens}[title={What you've built}]
One more word for telling what happened.
\end{grammarlens}

\subsection*{Your turn}
\begin{itemize}
\raggedright
  \item \emph{Say it:} \emph{tabāhī}
  \item \emph{Say it:} \emph{tabāhī}, once more
  \item \emph{Say it:} \emph{tūfān ne pūre śahar meṁ tabāhī kī.} — The storm wreaked havoc across the whole city.
  \item \emph{Recall:} say the Hindi for to flash, then the Hindi for torrential (of rain), then say \emph{tabāhī} again
\end{itemize}

\begin{tcolorbox}[breakable,colback=teal!4,colframe=teal!35!black,title={Before you move on}]
What does \hi{तबाही} mean? (\textquotedblleft{}Destruction; havoc\textquotedblright{}.) Say it once more.
\end{tcolorbox}
